\documentclass[10pt,a4paper]{amsart}
\usepackage[margin=19mm]{geometry}
\usepackage{amsmath,amssymb,mathtools}
\usepackage[scr=rsfs,cal=euler]{mathalfa}
\usepackage{xfrac}
\usepackage[unicode,hidelinks,bookmarks=false]{hyperref}
\newcommand{\ie}{i.e.\ }
\newcommand{\cf}{cf.\ }
\newcommand{\resp}{respectively\ }
\newcommand{\Subsection}{\subsection}
\providecommand{\nobreakdash}{\nobreak\hbox{-}}
\setlength{\emergencystretch}{2em}
\multlinegap0pt
\usepackage{xcolor}
\usepackage{amssymb}
\usepackage{float}
\usepackage{tikz}
\newtheorem{theorem}{Theorem}
\newtheorem{lemma}{Lemma}
\usepackage{cleveref}
\crefname{theorem}{Theorem}{Theorems}
\crefname{lemma}{Lemma}{Lemmas}
  \newcommand{\C}{\mathbb{C}} % complex numbers
  \DeclareMathOperator{\Diag}{Diag}
  \newcommand{\E}{\mathbb{E}}
  \DeclareMathOperator{\Est}{Est}
  \newcommand{\Exp}{\mathop{\mathbb{E}}}
  \newcommand{\F}{\mathbb{F}} % field
  \newcommand{\Lcal}{\mathcal{L}}
  \DeclareMathOperator{\Stab}{Stab}\DeclareMathOperator{\Spec}{Spec}
  \DeclareMathOperator{\U}{\mathcal{U}}
  \newcommand{\eps}{\varepsilon}
  \newcommand{\one}{\mathbf{1}}
  \newcommand{\pmset}[1]{\{-1,1\}^{#1}} % hypercube in +-1 basis
  \newcommand{\poly}{\mbox{\rm poly}}
  \DeclareMathOperator{\supp}{supp}
  \newcommand{\tp}{\mathsf{T}}
\title{An algorithmic Polynomial Freiman--Ruzsa theorem}
\author{Davi Castro-Silva, Jop Bri\"et, Srinivasan Arunachalam, Arkopal Dutt, Tom Gur}
\date{Source: arXiv:2604.04547v1 (CC BY 4.0)}
\begin{document}
\maketitle

\noindent\emph{Source and licence.} An algorithmic Polynomial Freiman--Ruzsa theorem. Version arXiv:2604.04547v1 (CC BY 4.0), distributed by arXiv under the Creative Commons Attribution 4.0 International licence, http://creativecommons.org/licenses/by/4.0/, as recorded in the arXiv licence metadata for this exact version and shown on the abstract page https://arxiv.org/abs/2604.04547v1. Full licence notice: \texttt{licenses/arxiv-2604.04547-CC-BY-4.0.txt}.\par\smallskip

\noindent\textit{Editorial scope.} Counted unit: the article's theorem thm:quadratic\_GL (source0.tex lines 424--429). The article's complete original proof of it is delivered with it (source0.tex lines 2211--2380, one \texttt{proof} environment of 170 lines, which the article places away from the statement). No mathematics is written, repaired or re-derived: the statement and the proof are the article's own lines, in the article's own notation, and no retained line is substituted or re-flowed.  Retained uncounted as the source-local context the proof needs: lines 1430--1440; lines 2149--2153.  Nothing was omitted from any retained span, so the package carries no omission record.  No retained line contains a citation, so the wrapper carries no bibliography and no bibliography entry.  No retained line contains a figure environment, an image-inclusion command or drawing code, so no image is delivered and the package declares no asset dependency.  Editorial formatting only, in addition: an AMS \texttt{amsart} preamble that restates the article's own declarations and macros used by the retained lines, namely the environments \texttt{theorem}, \texttt{lemma} on separate counters, together with the standard packages those lines need.  The retained environments print in this document as Theorem 1, Lemma 1 in order of appearance, whereas the article prints the same items with the numbering of its own full document.  The complete native source of the article as distributed in the arXiv e-print for this exact version is bundled unchanged as \texttt{sources/197889/source0.tex}.
\begin{lemma}[Description of stabilizer states] \label{lem:stab_explicit}
    Every stabilizer state $\phi \in \Stab(\F_2^n)$ can be written in the form
    \begin{equation} \label{eq:stab_explicit}
        \phi(x) = \alpha \sqrt{2^{n-\dim(V)}} \one_V(x+r) (-1)^{(x+r)^\tp A(x+r) + s\cdot(x+r)} i^{|d\circ (x+r)|},
    \end{equation}
    where $\alpha\in \U(1)$, $V\leq \F_2^n$ is a subspace, $A\in \F_2^{n\times n}$ is a matrix and $r, s , d \in \F_2^n$.
    Conversely, every function of the form~\eqref{eq:stab_explicit} is a stabilizer state, and its associated Lagrangian is
    \begin{equation} \label{eq:Lagrangian_explicit}
        \Lcal(\phi) = \big\{(h,\, Mh+w):\: h\in V,\, w\in V^\perp\big\} \quad \text{where} \quad M = A+A^\tp+\Diag(d).
    \end{equation}
\end{lemma}

\begin{theorem}[List-decoding stabilizer states] \label{thm:list}
    Let $f:\F_2^n \to \C$ be a 1-bounded function and let $\tau,\, \delta>0$ and $1/2<\gamma\leq 1$.
    There is a randomized algorithm that, when given query access to~$f$, returns a list of size $\log(\delta^{-1}) \big((\gamma-1/2) \tau\big)^{-O(\log(1/\tau))}$ which, with probability at least $1-\delta$, contains all $\gamma$-approximate local maximizers that have correlation at least~$\tau$ with~$f$.
    This algorithm makes $n^2 \log n\, \log(\delta^{-1}) \big((\gamma - \tfrac{1}{2}) \tau\big)^{-O(\log(1/\tau))}$ queries to~$f$ and has runtime $n^3 \log n\, \log(\delta^{-1}) \big((\gamma - \tfrac{1}{2}) \tau\big)^{-O(\log(1/\tau))}$.
\end{theorem}

\begin{theorem}[Quadratic Goldreich--Levin algorithm] \label{thm:quadratic_GL}
For every $\eps,\delta>0$, there exists a randomized algorithm such that the following holds.
Given query access to a function $f: \F_2^n \rightarrow \C$, with probability at least~$1 - \delta$, the algorithm outputs a quadratic polynomial $p: \F_2^n \rightarrow \F_2$ satisfying
$$\big|\Exp_{x \in \F_2^n} f(x) (-1)^{p(x)}\big| \geq \max_{q \text{ quadratic }} \big|\Exp_{x \in \F_2^n} f(x) (-1)^{q(x)}\big| - \eps.$$
This algorithm makes $n^2 \log n \log(1/\delta) (1/\varepsilon)^{O(\log(1/\varepsilon))}$ queries to $f$ and runs in time $n^3 \log n \log(1/\delta) (1/\varepsilon)^{O(\log(1/\varepsilon))}$.
\end{theorem}

\begin{proof}[ of \cref{thm:quadratic_GL}]
The main idea is to apply the algorithm from Theorem~\ref{thm:list} with suitably chosen parameters to obtain a bounded-size list containing all ``good'' stabilizer states, and then replace each of these good stabilizer states by a bounded number of (classical) quadratic phase functions.
Each such quadratic phase~$(-1)^q$ is obtained from its associated stabilizer state~$\phi$ by extending its support from (a coset of) a subspace~$V$ to the whole domain~$\F_2^n$.
We end the proof by showing that, with high probability, one of the quadratic phases thus obtained has almost-maximal correlation with~$f$;
by querying~$f$ a bounded number of times, we can estimate all of these correlations and pick up the highest one.

The full algorithm is given as follows:
\begin{enumerate}
    \item Apply the algorithm from Theorem~\ref{thm:list} with parameters $\tau = \eps$ and $\gamma = 1/2 + \eps^2$.
    We obtain a list~$L$ of size $\log(1/\delta) (1/\eps)^{O(\log(1/\eps))}$ which, with probability at least $1-\delta$, contains all stabilizer states that are $(1/2+\eps^2)$-approximate local maximizers of correlation for~$f$ and have correlation at least~$\eps$ with~$f$.
    
    \item Remove from~$L$ every stabilizer state whose support has codimension larger than $2\log(1/\eps)$.
    If~$L$ becomes empty after this step, end the algorithm and return the constant function $p \equiv 0$.
    Otherwise, initialize a list~$L'$ to be empty and continue the algorithm.
    
    \item For each stabilizer state $\phi\in L$, do the following:

    \noindent Write $\phi(x) = 2^{(n-d)/2} \one_{u+V}(x) (-1)^{q(x)} i^{|c\circ x|}$, where~$V$ is a subspace of dimension~$d$, $q: \F_2^n \to \F_2$ is a quadratic polynomial, and $u, c \in \F_2^n$ are vectors such that either $c=0$ or $c\notin V^\perp$.
    (This decomposition is possible due to \cref{lem:stab_explicit} and the remark following it.)
    \begin{itemize}
        \item[$(a)$] If $c=0$, add to~$L'$ the quadratic functions $x\mapsto q(x)+y\cdot x$ for all $y\in V^\perp$.

        \item[$(b)$] If $c\notin V^\perp$, let $U = \{c\}^\perp$ and let $v\in \F_2^n$ satisfy $c\cdot v = 1$, so that any $x\in\F_2^n$ has a unique representation of the form $x = z+bv$ for some $z\in U$ and $b\in \F_2$.
        By evaluating the map $x = z+bv \mapsto i^{|c\circ z| - 2|c\circ z\circ bv|}$ on all vectors $x\in \F_2^n$ with weight $|x| \leq 2$, find the polynomial $r\in \F_2[x_1,\dots,x_n]$ of degree at most~2 such that\footnote{Such a polynomial exists because the function $x = z+bv \mapsto i^{|c\circ z| - 2|c\circ z\circ bv|}$ is a non-classical quadratic phase function taking values in $\pmset{}$, and thus equals a classical quadratic phase $(-1)^{r(x)}$.}
        $(-1)^{r(z+bv)} = i^{|c\circ z| - 2|c\circ z\circ bv|}$.
        Add to~$L'$ the quadratic functions
        $$x\mapsto r(x)+q(x)+y\cdot x \quad \text{and} \quad x\mapsto r(x)+q(x)+ (y+c)\cdot x$$
        for all $y\in V^{\perp}$.
    \end{itemize}
    
    \item Query~$f$ at $m = \poly(\tfrac{1}{\eps} \log \tfrac{1}{\delta})$ randomly chosen points $x_1, \dots, x_m \in \F_2^n$ and compute
    $$\Est_q := \frac{1}{m} \sum_{j=1}^m f(x_j) (-1)^{q(x_j)}$$
    for all quadratic functions~$q$ in~$L'$.
    Output the one that attains the maximum value of~$|\Est_q|$.
\end{enumerate}

Note that, for each $\phi\in L$, the number of quadratic functions we add to~$L'$ at step~$(3)$ is at most~$2^{n-d+1}$.
Since $n-d \leq 2\log(1/\eps)$ because of step~$(2)$, it follows that the final list~$L'$ has size at most
$$2^{n-d+1} |L| \leq 2|L|/\eps^2 = \log(1/\delta) (1/\eps)^{O(\log(1/\eps))}.$$
In addition to the list-decoding subroutine from \cref{thm:list}, the most expensive step in this algorithm is step $(3)$, which takes $O(n^3 + n/\eps^2)$ time for each stabilizer state in the list~$L$.
The query and time complexities of the algorithm above thus match those stated in \cref{thm:quadratic_GL}.

Denote the (random) quadratic function output by this algorithm by~$p$.
We will show that, with probability at least~$1-2\delta$, this function satisfies
\begin{equation} \label{eq:high_corr}
    |\langle f,\, (-1)^{p(\cdot)}\rangle| > \|f\|_{u^3} - \eps,
\end{equation}
where $\|f\|_{u^3}$ (with a minuscule~$u$) denotes the maximum correlation $|\langle f,\, (-1)^{q(\cdot)}\rangle|$ of~$f$ with a quadratic polynomial $q\in \F_2[x_1, \dots, x_n]$.
This will complete the proof of the theorem.

Note that we may assume $\|f\|_{u^3} \geq \eps$, as otherwise any quadratic function will satisfy equation~$\eqref{eq:high_corr}$.
We can also assume that~$\eps \leq 1/100$, which will allow us to bound certain expressions more easily.
The heart of the argument is given in the following result:

\begin{lemma} \label{lem:find_quad}
    Assume that~$\eps \leq 1/100$ and~$\|f\|_{u^3} \geq \eps$.
    Then, with probability at least~$1-\delta$, there exists a quadratic function~$q$ in~$L'$ satisfying
    $$|\langle f,\, (-1)^{q(\cdot)}\rangle| \geq \|f\|_{u^3} - \eps/2.$$
\end{lemma}

\begin{proof}
Let $p^*: \F_2^n \to \F_2$ be a quadratic function attaining maximum correlation with~$f$:
$$\big|\E_{x\in \F_2^n} f(x) (-1)^{p^*(x)}\big| = \|f\|_{u^3}.$$
Consider the stabilizer state $\phi_0 := (-1)^{p^*(\cdot)}$, and denote $\gamma = 1/2+\eps^2$.
If~$\phi_0$ is a $\gamma$-approximate local maximizer of correlation with~$f$, then with probability at least~$1-\delta$ it will appear in the list~$L$
(and thus~$p^*$ will appear in~$L'$).

Now suppose~$\phi_0$ is not a $\gamma$-approximate local maximizer of correlation with~$f$.
There must then exist a ``neighbor'' stabilizer state~$\phi_1$ satisfying
$$|\langle\phi_0,\, \phi_1\rangle|^2 = 1/2 \quad \text{and} \quad |\langle f,\, \phi_1 \rangle|^2 > \gamma^{-1} |\langle f,\, \phi_0 \rangle|^2.$$
If~$\phi_1$ is a $\gamma$-approximate local maximizer, then it will appear in the list~$L$ with probability at least~$1-\delta$.
Otherwise, we can keep choosing stabilizer states~$\phi_{i+1}$ satisfying
$$|\langle\phi_i,\, \phi_{i+1}\rangle|^2 = 1/2 \quad \text{and} \quad |\langle f,\, \phi_{i+1} \rangle|^2 > \gamma^{-1} |\langle f,\, \phi_i \rangle|^2$$
until at last we arrive at some~$\phi_t$ which is a $\gamma$-approximate local maximizer of correlation with~$f$.
This must stop at some point because $|\langle f,\, \phi_0 \rangle| = \|f\|_{u^3} \geq \eps$ and we always have $|\langle f,\, \phi_i \rangle| \leq \|f\|_2$ by Cauchy-Schwarz.
The final stabilizer state~$\phi_t$ will then appear in list~$L$ with probability at least~$1-\delta$, and it satisfies
\begin{equation} \label{eq:corr_phi_t}
    |\langle f,\, \phi_t \rangle|^2 > \gamma^{-t} |\langle f,\, \phi_0 \rangle|^2 = (1/2+\eps^2)^{-t} \|f\|_{u^3}^2.
\end{equation}

Let us write
$$\phi_t(x) = 2^{(n-d)/2} \one_{u+V}(x) (-1)^{q^*(x)} i^{|c\circ x|},$$
where~$V$ is a subspace of dimension $d$, $q^*: \F_2^n \to \F_2$ is a quadratic polynomial and $u, c \in \F_2^n$ are vectors such that either~$c=0$ or $c\notin V^{\perp}$.
Since
$$\eps \leq |\langle f,\, \phi_0 \rangle| \leq |\langle f,\, \phi_t \rangle| \leq 2^{(n-d)/2} \E_{x\in \F_2^n} \one_{u+V}(x) |f(x)| = 2^{-(n-d)/2},$$
we conclude that $n-d \leq 2\log(1/\eps)$, and so~$\phi_t$ will not get removed in step~$(2)$ of the~algorithm.

Next, we relate the dimension~$d$ of~$V$ with the number~$t$ of steps we took until we arrived at~$\phi_t$.
For each $0\leq i\leq t$, denote by $\dim(\phi_i)$ the dimension of the subspace on which the $i$-th stabilizer state~$\phi_i$ is supported.
Since $|\langle\phi_i,\, \phi_{i+1}\rangle|^2 = 1/2$ while $|\phi_j(\cdot)| = 2^{(n-\dim(\phi_j))/2} \one_{\supp(\phi_j)}(\cdot)$, we conclude that
$\dim(\phi_{i+1}) \geq \dim(\phi_i)-1$.
Moreover, in the case where $\dim(\phi_{i+1}) = \dim(\phi_i)-1$, the two stabilizer states~$\phi_i$ and~$\phi_{i+1}$ must be proportional to one another inside the support of~$\phi_{i+1}$.
As~$\phi_0 = (-1)^{p^*(\cdot)}$ while~$\phi_t$ has a nontrivial non-classical component~$i^{|c\circ x|}$ if $c\notin V^{\perp}$, it follows that $\dim(\phi_t) \geq \dim(\phi_0) - t + \one_{c\notin V^{\perp}}$.
We conclude that $t \geq n-d + \one_{c\notin V^{\perp}}$.

Using the fact that $\E_{y\in V^{\perp}} (-1)^{y\cdot x} = \one_V(x)$, we see that
\begin{align*}
    |\langle f,\, \phi_t \rangle|^2
    &= 2^{n-d} \big|\E_{x\in \F_2^n} f(x) \one_{V}(x+u) (-1)^{q^*(x)} i^{-|c\circ x|}\big|^2 \\
    &= 2^{n-d} \big|\E_{y\in V^{\perp}} \E_{x\in \F_2^n} f(x) (-1)^{y\cdot (x+u)} (-1)^{q^*(x)} i^{-|c\circ x|}\big|^2 \\
    &\leq 2^{n-d} \E_{y\in V^{\perp}} \big|\E_{x\in \F_2^n} f(x) (-1)^{q^*(x) + y\cdot x} i^{-|c\circ x|}\big|^2,
\end{align*}
and thus there exists some $y^* \in V^{\perp}$ such that
\begin{equation} \label{eq:ystar}
    \big|\E_{x\in \F_2^n} f(x) (-1)^{q^*(x) + y^* \cdot x} i^{-|c\circ x|}\big|^2 \geq 2^{-(n-d)} |\langle f,\, \phi_t \rangle|^2.
\end{equation}
Recall that, if $\phi_t \in L$ (which happens with probability at least~$1-\delta$), then the quadratic functions $x\mapsto r(x) + q^*(x) + y^*\cdot x$ and $x\mapsto r(x) + q^*(x) + (y^*+c)\cdot x$ will both be in~$L'$
(we will recall the definition of $r\in \F_2[x_1, \dots, x_n]$ below).
It then suffices to show that one of these functions has correlation at least $\|f\|_{u^3}-\eps/2$ with~$f$.

We separate the proof into two cases:
$c=0$ and $c\notin V^{\perp}$.
If~$c=0$, then $r\equiv 0$, and we obtain from combining~\eqref{eq:corr_phi_t} and~\eqref{eq:ystar} that
$$\big|\E_{x\in \F_2^n} f(x) (-1)^{q^*(x) + y^* \cdot x}\big|^2 \geq 2^{-(n-d)} (1/2+\eps^2)^{-t} \|f\|_{u^3}^2.$$
Using that $n-d \leq 2\log(1/\eps)$ and $t\geq n-d$, we conclude that
$$2^{-(n-d)} (1/2+\eps^2)^{-t} \geq 2^{-(n-d)} (1/2+\eps^2)^{-(n-d)} \geq (1+2\eps^2)^{-2\log(1/\eps)}.$$
This last expression is at least $(1-\eps/2)^2$ when $\eps \leq 1/100$, which implies that
$$\big|\E_{x\in \F_2^n} f(x) (-1)^{q^*(x) + y^* \cdot x}\big| \geq \|f\|_{u^3} - \eps/2$$
as wished.

In the case where $c\notin V^{\perp}$, let~$U \leq \F_2^n$ be the subspace orthogonal to~$c$, and let~$v \in \F_2^n\setminus U$.
Any $x\in \F_2^n$ can be written in a unique way as $x = z+bv$, where $z\in U$ and $b\in \F_2$.
Note that $|c\circ (z+bv)| = |c\circ z| + |c\circ bv| - 2|c\circ z\circ bv|$.
Define $h:\F_2^n\to \C$ by
$$h(z+bv) = i^{|c\circ z|  - 2|c\circ z\circ bv|} \quad \text{where $z\in U$, $b\in \F_2$.}$$
Then~$h$ is a non-classical quadratic phase function taking values in $\pmset{}$, which implies it must be classical:
there exists a polynomial $r\in \F_2[x_1,\dots,x_n]$ of degree at most~2 such that $h(x) = (-1)^{r(x)}$ for all~$x\in \F_2^n$.

For any function $g: \F_2^n \to \C$, we have
\begin{align*}
    \big|\E_{x\in \F_2^n} g(x) i^{-|c\circ x|}\big|^2
    &= \big|\E_{b\in \F_2} \E_{z\in U} g(z+bv) i^{-|c\circ (z+bv)|} \big|^2 \\
    &= \big|\E_{b\in \F_2} \E_{z\in U} g(z+bv) (-1)^{r(z+bv)} i^{-|c\circ bv|} \big|^2.
\end{align*}
By the Cauchy-Schwarz inequality, this expression is at most
\begin{equation*}
    \E_{b\in \F_2}\big|\E_{z\in U} g(z+bv)(-1)^{r(z+bv)}\big|^2.
\end{equation*}
By Parseval's identity on~$\F_2$ we get
\begin{align*}
    \E_{b\in \F_2} |\E_{z\in U} g(z+bv) (-1)^{r(z+bv)}|^2
    &= \sum_{a\in \F_2} |\E_{b\in \F_2} \E_{z\in U} g(z+bv) (-1)^{r(z+bv) + ab}|^2 \\
    &= \sum_{a\in \F_2} |\E_{b\in \F_2} \E_{z\in U} g(z+bv) (-1)^{r(z+bv) + a c\cdot (z+bv)}|^2 \\
    &= |\E_{x\in \F_2^n} g(x)(-1)^{r(x)}|^2 + |\E_{x\in \F_2^n} g(x) (-1)^{r(x) + c\cdot x}|^2,
\end{align*}
from which we conclude that
$$|\E_{x\in \F_2^n} g(x)(-1)^{r(x)}|^2 + |\E_{x\in \F_2^n} g(x) (-1)^{r(x) + c\cdot x}|^2 \geq \big|\E_{x\in \F_2^n} g(x) i^{-|c\circ x|}\big|^2.$$

Using this last inequality for the function $g(x) = f(x) (-1)^{q^*(x) + y^*\cdot x}$, we obtain
\begin{align*}
    \max\big\{|\E_{x\in \F_2^n} f(x) (-1)^{r(x) + q^*(x) + y^* \cdot x}|^2,\: &|\E_{x\in \F_2^n} f(x) (-1)^{r(x)+q^*(x) + (y^*+c) \cdot x}|^2 \big\} \\
    &\geq \frac{1}{2} \big|\E_{x\in \F_2^n} f(x) (-1)^{q^*(x) + y^* \cdot x} i^{-|c\circ x|}\big|^2 \\
    &\geq \frac{1}{2^{n-d+1}} |\langle f,\, \phi_t \rangle|^2 \\
    &\geq \frac{1}{2^{n-d+1}} \Big(\frac{1}{1/2+\eps^2}\Big)^t \|f\|_{u^3}^2,
\end{align*}
where we used inequalities~\eqref{eq:ystar} and~\eqref{eq:corr_phi_t} respectively.
Since in this case we have $t\geq n-d+1$ and $n-d\leq 2\log(1/\eps)$, the last expression is at least
$$\Big(\frac{1}{1 + 2\eps^2}\Big)^{2\log(1/\eps)+1} \|f\|_{u^3}^2 \geq (1-\eps/2)^2 \|f\|_{u^3}^2,$$
where we use that~$\eps \leq 1/100$. 
This concludes the proof of the lemma.
\end{proof}

Using our bound on the size of the list~$L'$ and the Chernoff bound we conclude that, with probability at least~$1-\delta$, we have
$$\big|\Est_q - \langle f,\, (-1)^{q(\cdot)}\rangle\big| < \eps/4 \quad \text{for all $q\in L'$.}$$
Recall that we denote by~$p$ the random polynomial output by the algorithm.
If the above estimate holds, we get
$$|\langle f,\, (-1)^{p(\cdot)}\rangle| > |\Est_p| - \frac{\eps}{4} = \max_{q \in L'} |\Est_q| - \frac{\eps}{4} > \max_{q \in L'} |\langle f,\, (-1)^{q(\cdot)}\rangle| - \frac{\eps}{2}.$$
By Lemma~\ref{lem:find_quad}, we have that $\max_{q \in L'} |\langle f,\, (-1)^{q(\cdot)}\rangle| \geq \|f\|_{u^3} - \eps/2$ with probability at least $1-\delta$.
This implies that inequality~\eqref{eq:high_corr} holds with probability at least $1-2\delta$, finishing the proof of the theorem.
\end{proof}

\end{document}
